\documentclass[10pt,a4paper]{amsart}
\usepackage{amsmath,amssymb,amsthm,graphicx}
\usepackage{verbatim}
\newenvironment{reference}{\comment}{\endcomment}
%\newenvironment{reference}{}{}
\newenvironment{slogan}{\comment}{\endcomment}
\newenvironment{history}{\comment}{\endcomment}

% For commutative diagrams we use Xy-pic
\usepackage[all]{xy}

% We use 2cell for 2-commutative diagrams.
\xyoption{2cell}
% Theorem environments.
%
\theoremstyle{plain}
\newtheorem{theorem}[subsection]{Theorem}
\newtheorem{proposition}[subsection]{Proposition}
\newtheorem{lemma}[subsection]{Lemma}

\theoremstyle{definition}
\newtheorem{definition}[subsection]{Definition}
\newtheorem{example}[subsection]{Example}
\newtheorem{exercise}[subsection]{Exercise}
\newtheorem{situation}[subsection]{Situation}

\theoremstyle{remark}
\newtheorem{remark}[subsection]{Remark}
\newtheorem{remarks}[subsection]{Remarks}

\numberwithin{equation}{subsection}

% Macros
%
\def\lim{\mathop{\mathrm{lim}}\nolimits}
\def\colim{\mathop{\mathrm{colim}}\nolimits}
\def\Spec{\mathop{\mathrm{Spec}}}
\def\Hom{\mathop{\mathrm{Hom}}\nolimits}
\def\Ext{\mathop{\mathrm{Ext}}\nolimits}
\def\SheafHom{\mathop{\mathcal{H}\!\mathit{om}}\nolimits}
\def\SheafExt{\mathop{\mathcal{E}\!\mathit{xt}}\nolimits}
\def\Sch{\mathit{Sch}}
\def\Mor{\mathop{\mathrm{Mor}}\nolimits}
\def\Ob{\mathop{\mathrm{Ob}}\nolimits}
\def\Sh{\mathop{\mathit{Sh}}\nolimits}
\def\NL{\mathop{N\!L}\nolimits}
\def\CH{\mathop{\mathrm{CH}}\nolimits}
\def\proetale{{pro\text{-}\acute{e}tale}}
\def\etale{{\acute{e}tale}}
\def\QCoh{\mathit{QCoh}}
\def\Ker{\mathop{\mathrm{Ker}}}
\def\Im{\mathop{\mathrm{Im}}}
\def\Coker{\mathop{\mathrm{Coker}}}
\def\Coim{\mathop{\mathrm{Coim}}}

% Boxtimes
%
\DeclareMathSymbol{\boxtimes}{\mathbin}{AMSa}{"02}

%
% Macros for moduli stacks/spaces
%
\def\QCohstack{\mathcal{QC}\!\mathit{oh}}
\def\Cohstack{\mathcal{C}\!\mathit{oh}}
\def\Spacesstack{\mathcal{S}\!\mathit{paces}}
\def\Quotfunctor{\mathrm{Quot}}
\def\Hilbfunctor{\mathrm{Hilb}}
\def\Curvesstack{\mathcal{C}\!\mathit{urves}}
\def\Polarizedstack{\mathcal{P}\!\mathit{olarized}}
\def\Complexesstack{\mathcal{C}\!\mathit{omplexes}}
% \Pic is the operator that assigns to X its picard group, usage \Pic(X)
% \Picardstack_{X/B} denotes the Picard stack of X over B
% \Picardfunctor_{X/B} denotes the Picard functor of X over B
\def\Pic{\mathop{\mathrm{Pic}}\nolimits}
\def\Picardstack{\mathcal{P}\!\mathit{ic}}
\def\Picardfunctor{\mathrm{Pic}}
\def\Deformationcategory{\mathcal{D}\!\mathit{ef}}
\usepackage{cleveref}
\crefname{originaltheorem}{original theorem}{original theorems}
\Crefname{originaltheorem}{Original theorem}{Original theorems}
\makeatletter
\newlabel{algebra-theorem-chevalley}{{00FE}{source}{}{sourcepointer-1}{}}
\newlabel{algebra-theorem-chevalley@cref}{{[originaltheorem][0][]00FE}{[1][0][]source}{}{}{}}
\makeatother
\crefname{originalsection}{original section}{original sections}
\Crefname{originalsection}{Original section}{Original sections}
\makeatletter
\newlabel{morphisms-section-constructible}{{054H}{source}{}{sourcepointer-2}{}}
\newlabel{morphisms-section-constructible@cref}{{[originalsection][0][]054H}{[1][0][]source}{}{}{}}
\makeatother
\title{A finitely presented map of affine schemes taking a locally closed subset to a set that is not a finite union of locally closed subsets}
\author{The Stacks Project Authors}
\date{stacks-project:examples.tex:local-snapshot-20261012; GNU FDL 1.2 or later}
\begin{document}
\maketitle
\noindent\textit{Editorial note.} All delivered mathematical text is the author's own, copied from the pinned Stacks Project chapter file examples.tex; no statement or proof step was rewritten, reordered or removed.\par\medskip
\noindent\textit{Editorial note.} The counted proof is the author's complete argument: The counted proof adds no mathematics of its own; the shift-self-identification of B with B cap (y\_1 = 0) is the pivot of the context induction.\par\medskip
\noindent\textit{Editorial note.} The label lemma-no-chevalley names the counted statement, and the proof environment immediately following it in the source is delivered verbatim as its proof.\par\medskip
\noindent\textit{Editorial note.} The counted statement is carried by the author's own lemma environment in the chapter source.\par\medskip
\noindent\textit{Allowed standard prerequisites (editorial).}\begin{itemize}
\item Context span 2112-2164: over a field k of characteristic 0 the projection f: X = Spec(k[t, x\_1, x\_2, ..., y\_1, y\_2, ...]) -> Spec(k[x\_1, x\_2, ..., y\_1, y\_2, ...]) = Y is of finite presentation; Z is the closed subset cut out by x\_1 (t - 1) = 0, x\_2 (t - 1)(t - 2) = 0, x\_3 (t - 1)(t - 2)(t - 3) = 0, ...; U = union of the opens U\_j where y\_j (t - 1)...(t - j) is nonzero; on each U\_j the image f(Z cap U\_j) is the set of points where x\_1, ..., x\_j are zero and y\_j is nonzero
\item Context claim and proof: B = f(Z cap U) = union\_j f(Z cap U\_j) is not a finite union of locally closed subsets of Y; the induction on n shows any cover B = A\_1 union ... union A\_m has m >= n for all n (the closure of B is (x\_1 = 0), some A\_i is open in (x\_1 = 0) but cannot contain a point with y\_1 = 0 since for points of B y\_1 = 0 forces x\_2 = 0; the right shift map x\_i -> x\_\{i+1\}, y\_i -> y\_\{i+1\} identifies B with B cap (y\_1 = 0))
\item Chevalley's theorem (Algebra, theorem-chevalley; Morphisms, section-constructible) gives constructibility of images of constructible sets under finitely presented morphisms - the positive backdrop stated in the section preamble
\end{itemize}\medskip
\ifdefined\hypertarget\hypertarget{sourcepointer-1}{}\fi
\noindent\textit{Source reference.} Original theorem 00FE: sources/stacks-chow/tags.txt, source lines 568--568. Prerequisite pointer invoked at 2179; it names a standard fact of a sibling chapter that the counted argument uses, and every step of the counted proof built on it is retained verbatim in the delivered body.\par\medskip
\ifdefined\hypertarget\hypertarget{sourcepointer-2}{}\fi
\noindent\textit{Source reference.} Original section 054H: sources/stacks-chow/tags.txt, source lines 6640--6640. Prerequisite pointer invoked at 2180; it names a standard fact of a sibling chapter that the counted argument uses, and every step of the counted proof built on it is retained verbatim in the delivered body.\par\medskip
\noindent
Chevalley's theorem says that the image of a constructible set
by a finitely presented morphism of affine schemes is constructible, see
Algebra, Theorem \ref{algebra-theorem-chevalley} and
Morphisms, Section \ref{morphisms-section-constructible}.
We will see the same thing does not hold
for images of locally closed subsets.

\medskip\noindent
Let $k$ be a field of characteristic $0$. Consider the projection morphism
$$
f :
X = \Spec(k[t, x_1, x_2, \ldots, y_1, y_2, \ldots])
\longrightarrow
\Spec(k[x_1, x_2, \ldots, y_1, y_2, \ldots]) = Y
$$
This is a morphism of finite presentation.
Let $Z$ be the closed subset of $X$ defined by
$$
x_1(t - 1) = 0,\quad
x_2(t - 1)(t - 2) = 0,\quad
x_3(t - 1)(t - 2)(t - 3) = 0,\quad \ldots
$$
Let $U = \bigcup_{j \geq 1} U_j$ be the open of $X$ defined by
$$
U_j = \text{points where }y_j(t - 1)(t - 2) ... (t - j)\text{ is nonzero}
$$
Then we have
$$
f(Z \cap U_j) =
\text{points where }x_1, \ldots, x_j\text{ are zero and }y_j\text{ is nonzero}
$$
We claim that $B = f(Z \cap U) = \bigcup_{j \geq 1} f(Z \cap U_j)$ is
not a finite union of locally closed subsets of $Y$.

\medskip\noindent
Proof of the claim. Suppose that $B = A_1 \cup \cdots \cup A_m$
is a finite cover of $B$ by locally closed subsets of $Y$.
We will show by induction on $n$ that $m \geq n$. The base
case $n = 1$ is OK as $B$ is nonempty. Assume $n > 1$ and that
the induction hypothesis holds for $n - 1$. Since the closure of $B$ is
$(x_1 = 0)$, one of the $A_i$ must contain some nonempty open subset of
$(x_1 = 0)$. Then $A_i$ must be open in $(x_1 = 0)$.
But any such open subset cannot contain a point with
$y_1 = 0$; indeed, for points of $B$, $y_1 = 0$ forces $x_2 = 0$,
and this shows $B$ contains no neighborhood of $(x, y)$ inside $(x_1 = 0)$.
Therefore, the remaining $m - 1$ elements restrict
to a constructible cover of $B \cap (y_1 = 0)$.
However, observe that the right shift map
$x_i \mapsto x_{i + 1}$, $y_i \mapsto y_{i + 1}$ identifies $B$
with $B \cap (y_1 = 0)$! Thus by induction hypothesis, we see that
$m - 1 \geq n - 1$ and we conclde $m \geq n$. This finishes
the proof of the induction step and thereby establishes the claim.

\begin{lemma}
\label{lemma-no-chevalley}
There exists a morphism $f : X \to Y$ of finite presentation
between affine schemes and a locally closed subset $T$ of $X$
such that $f(T)$ is not a finite union of locally closed subsets of $Y$.
\end{lemma}

\begin{proof}
See discussion above.
\end{proof}

\end{document}
